\chapter{Pendahuluan}\label{bab:pendahuluan}

\section{Latar Belakang}\label{sec:latar-belakang}

Dokumen hasil pindai menyimpan informasi pada dua lapis sekaligus. Lapis pertama adalah
\textbf{tata letak}: posisi blok, kerapatan tulisan, tabel, kop surat, dan tanda tangan.
Lapis kedua adalah \textbf{isi tekstual}: kata-kata seperti \textit{Dear Sir},
\textit{Memorandum}, atau \textit{Curriculum Vitae} yang langsung menandai jenis dokumen.
Lapis pertama terlihat langsung pada citra. Lapis kedua tersembunyi di dalam piksel dan
harus dikeluarkan lebih dulu melalui \textit{optical character recognition} (OCR).
Karena teksnya diturunkan dari citra, \textbf{mutu modalitas teks dibatasi oleh mutu
OCR-nya}.

Pada laporan Progress~1, kelompok kami mengusulkan replikasi model dua cabang
\textcite{audebert2020}. Model ini menggabungkan MobileNetV2 untuk citra halaman dengan
FastText dan CNN~1D untuk teks hasil Tesseract. Replikasi diuji pada dataset Tobacco3482,
dan Progress~1 mengajukan tiga pertanyaan:
\begin{enumerate*}[label=(\alph*)]
  \item seberapa besar akurasi naik bila kedua modalitas digabung;
  \item strategi penggabungan mana yang bekerja; dan
  \item apakah model gabungan sungguh memakai kedua sumber.
\end{enumerate*}

Sejak Progress~1, replikasi tersebut telah dijalankan penuh, dan ketiga pertanyaan itu
telah terjawab (Subbab~\ref{sec:studi-pendahuluan}). Jawaban pertanyaan (c) yang paling
menentukan arah penelitian ini. Secara angka, replikasi berhasil: model gabungan
mencapai akurasi 0,834, di atas baseline citra (0,809) dan baseline teks (0,727). Namun
ketika seluruh masukan teks dinolkan saat inferensi, akurasinya hanya turun
\textbf{0,0015}. Model gabungan itu ternyata \textbf{diam-diam unimodal}: kenaikan
akurasinya tidak berasal dari teks. Padahal komplementaritas kedua modalitas jelas ada,
karena 8,9\% dokumen uji hanya dapat diklasifikasikan dengan benar melalui teks.

Temuan ini menyingkap kelemahan cara bidang ini melaporkan hasil. Studi fusion teks--citra
umumnya hanya melaporkan akurasi gabungan. Padahal, seperti ditunjukkan replikasi kami,
akurasi gabungan yang tinggi tidak membuktikan bahwa kedua modalitas benar-benar dipakai.

Studi terkini mewarisi kelemahan yang sama dan menambah satu kelemahan lain.
\textcite{kashyap2026} menggabungkan EfficientNetB3 dan BERT, dan melaporkan akurasi
92,00\% pada Tobacco3482. Akan tetapi, mesin OCR yang menghasilkan teksnya hanya disebut
\enquote{mesin OCR komersial}, tanpa nama. Akibatnya, keunggulan cabang teks mereka tidak
dapat dipisahkan dari mesin OCR yang tidak diketahui itu.

Di tingkat bidang, meta-analisis atas 139 studi oleh \textcite{mironczuk2026} menemukan
bahwa fusion multimodal menaikkan akurasi rata-rata 5,28 poin. Namun meta-analisis itu
\emph{tidak dapat} mengontrol kualitas OCR, karena studi primer jarang melaporkannya.
Hanya 11,8\% studi multimodal yang memakai uji statistik. Perbandingan mesin OCR yang ada,
misalnya \textcite{francis2025}, mengukur ketepatan OCR itu sendiri, bukan dampaknya pada
tugas lanjutan. Padahal galat OCR terbukti merambat ke tugas berikutnya
\parencite{zhang2025}.

Tanggapan dosen pengampu atas Progress~1 mengarahkan penelitian ini ke celah tersebut.
Arahan itu meliputi eksplorasi mesin OCR selain Tesseract (EasyOCR, PaddleOCR,
Keras-OCR), eksplorasi pada sisi embedding, perumusan celah penelitian yang eksplisit,
dan referensi utama yang terbit paling awal tahun 2025. Riwayat lengkap pemilihan
referensi dan tanggapan dosen dicatat pada Lampiran~\ref{lam:riwayat}.

Penelitian ini karena itu menjadikan \textbf{mesin OCR} dan \textbf{embedding teks}
sebagai variabel. Pengaruh keduanya diukur bukan hanya terhadap akurasi, tetapi juga
terhadap \emph{seberapa jauh model gabungan benar-benar memakai} modalitas teks.

\section{Rumusan Masalah}\label{sec:rumusan-masalah}

\begin{enumerate}[label=RM\arabic*.,leftmargin=1.2cm,labelwidth=1.2cm]
  \item Seberapa besar pilihan mesin OCR (Tesseract, EasyOCR, PaddleOCR) mengubah akurasi
  cabang teks pada klasifikasi dokumen Tobacco3482?
  \item Apakah embedding kontekstual (BERT) memperkecil selisih antar-mesin OCR
  dibandingkan embedding subword statis (FastText)? Dengan kata lain, adakah interaksi
  antara mesin OCR dan jenis embedding?
  \item Apakah cabang teks yang lebih kuat membuat model gabungan lebih
  \emph{memakai} modalitas teks? Apakah aturan fusion rata-rata logit lebih tahan
  terhadap cabang yang timpang dibandingkan fusion tingkat fitur?
\end{enumerate}

\section{Tujuan Penelitian}\label{sec:tujuan}

\begin{enumerate}[label=T\arabic*.,leftmargin=1.2cm,labelwidth=1.2cm]
  \item Mengukur pengaruh mesin OCR terhadap akurasi dan macro F1 cabang teks, disertai
  uji statistik antar-mesin.
  \item Mengukur interaksi mesin OCR dan jenis embedding melalui desain faktorial.
  \item Mengukur pemanfaatan modalitas teks pada model gabungan melalui uji
  \textit{missing modality}, serta membandingkan aturan rata-rata logit
  \parencite{kashyap2026} dengan fusion tingkat fitur.
\end{enumerate}

\section{Manfaat Penelitian}\label{sec:manfaat}

\begin{enumerate}
  \item \textbf{Manfaat akademik.} Penelitian ini menyediakan protokol evaluasi fusion
  yang melampaui akurasi gabungan: uji \textit{missing modality}, beberapa seed, dan uji
  statistik. Protokol ini menjawab kekurangan pelaporan yang diidentifikasi
  \textcite{mironczuk2026}.
  \item \textbf{Manfaat praktis.} Penelitian ini memberi dasar empiris untuk memilih mesin
  OCR open-source pada sistem klasifikasi dokumen. Dasar ini diukur dari dampaknya pada
  klasifikasi, bukan dari ketepatan OCR semata.
  \item \textbf{Bahan tahap lanjutan.} Kotak posisi teks dari setiap mesin OCR disimpan
  sejak awal sebagai bahan tahap berikutnya, yaitu teks sebagai objek
  (Subbab~\ref{sec:lanjutan}).
\end{enumerate}

\section{Batasan Masalah}\label{sec:batasan}

\begin{enumerate}
  \item Dataset yang dipakai hanya Tobacco3482 (3.482 dokumen, 10 kelas). RVL-CDIP
  \parencite{harley2015} tidak dipakai karena ukurannya ($\pm$400.000 dokumen) di luar kemampuan perangkat yang
  tersedia. Akibatnya, angka penelitian ini \textbf{tidak sebanding langsung} dengan
  \textcite{kashyap2026}, yang melatih modelnya pada RVL-CDIP.
  \item Kualitas OCR diukur melalui proksi. Tobacco3482 tidak memiliki transkripsi acuan,
  sehingga \textit{character error rate} tidak dapat dihitung.
  \item Cabang citra dibuat tetap (MobileNetV2) di semua kondisi, supaya perbedaan hasil
  hanya berasal dari sisi teks.
  \item Seluruh komputasi dijalankan pada satu laptop dengan GPU 4~GB VRAM.
\end{enumerate}
